\DTMsavedate{mydate}{2022-06-30}
\maketitle{Digital Devices and Systems}{Electronic Engineering}{\DTMusedate{mydate}}

% ----------------------------------------------------- %
% COURSE INFO
\subsection*{Course Information}
\vspace{0ex}%
\noindent\parbox{0.5\textwidth}{%
    \noindent\begin{tabular}{@{}ll}
                 \textsf{Course:}          & Digital Devices and Systems \\
                 \textsf{Course Code:}         & 	TEE 2211    \\
                 \textsf{Credit Hours:}    & 	10
    \end{tabular}}
\vspace{2ex}\hrule\vspace{2ex}

% ----------------------------------------------------- %
% INSTRUCTOR INFO
\subsection*{Lecturer Information}
\noindent\parbox{0.5\textwidth}{%
    \noindent\begin{tabular}{@{}ll}
                 \textsf{Name:}         &    Svetlana A. Bebova           \\
% \textsf{Phone:}        &    \phone          \\
\textsf{Email:}        &    \href{mailto:swetlana.bebova@nust.ac.zw}{swetlana.bebova@nust.ac.zw} \\
\textsf{Qualifications:}        &  MSc in Radio Electronic Engineering (Technical University of Sofia)
    \end{tabular}}
\vspace{2ex}\hrule\vspace{2ex}


% ----------------------------------------------------- %
% COURSE DESCRIPTION
\subsection*{Course Synopsis}
The module focuses on Flip-Flops review i.e.\ master-slave Flip-Flops, Shift registers.
Counters i.e asynchronous with mod(numbers) < 2N, synchronous, down counters, up/down counters, integrated circuits counters.
Memory devices i.e.\ registers, RAM, ROM and PROM.
Combinational circuits i.e.\ adders, multiplexer, decoders and code coverters.
Arithmetic circuits, arithmetic logic unit.
Design of combinational circuits using MSI devices.
Sequential system designs.
Design of synchronous sequential systems.
\vspace{2ex} \hrule\vspace{2ex}
% ----------------------------------------------------- %
% COURSE TOPICS
\subsection*{Topics}
\begin{outline}
    \1 Memory devices
    \1 ROM
    \1 RAM
    \1 PROM
    \1 PLA
    \1 Combinational  circuits
    \begin{itemize}
        \item Adders
        \item Multiplexers
        \item Decoders
        \item Code converters
    \end{itemize}
    \1 Design of combinational  circuits using MSI devices
    \1 Arithmetic Logic Unit (ALU)
    \1 Sequential system designs
    \1 Design of  asynchronous sequential systems
\end{outline}
\vspace{2ex} \hrule\vspace{2ex}

% ----------------------------------------------------- %
% Students Assenssment
\subsection*{Assessment of Students}
\begin{outline}
    \1 Students are assessed through written assignments and written in-class tests that shall form part of course work assessment mark.
    \1 There shall be a final examination at the end of the semester.
\end{outline}
\vspace{2ex}\hrule\vspace{2ex}

% ----------------------------------------------------- %
% Grade Evaluation
\subsection*{Course Evaluation and Grading Policies}
Grades are based on:
Continuous assessment (\qty{25}{\percent}),
Final Exam (\qty{75}{\percent})
\vspace{2ex}\hrule\vspace{2ex}

% ----------------------------------------------------- %
% EQUIPMENT
% https://sites.udel.edu/computing-purchases/personal-specs/
% \subsection*{Essential Equipment}

% \vspace{2ex}\hrule\vspace{2ex}

% Policies and Other information
\subsection*{Policies and Other Information}
% Conduct
\subsection*{Key Text}
\begin{itemize}
    \item Digital systems: Principles and applications, 1985 by  R. J.Tocci
    \item Digital Fundamentals, 1990 by  T. L. Floyd
    \item Logic and computer design fundamentals, 1997 by  M. M. Mano, C. R.Kime
    \item Electronic logic systems, 1994 by A.E.A. Almain
    \item Digital logic and computer design, 1979 by M. M. Mano
    \item Introduction to Digital Logic Design, 1994 by J. P. Hayes
\end{itemize}
\vspace{2ex}\hrule
\vspace{2ex}\hrule\vspace{2ex}